\chapter{Conclusion}
\label{ch:conclusion}

\section{What was promised}

The proposal committed to an instrument. It said so three times: that the primary deliverable
is a functional Asynchronous Measurement Apparatus, that this is ``not a proposed future
application and not a conceptual description'', and that the project ``proposes a thermometer,
and delivers it as working, open-source code that can be run today''.

That framing governs how this work should be judged. A thermometer that reads a value nobody
expected is still a delivered thermometer. Failing to ship it would have been the failure.

\section{What was delivered}

\subsection{The apparatus}

The instrument exists and runs. Given a passage of text it synthesises speech, obtains word
timings, computes text and audio embeddings, passes them to the released encoder, and returns
a predicted response over $20{,}484$ cortical surface points, from which the index is
computed. It is open source, it is deployed, and it was run end to end over 400 items.

Its outputs are reproducible: the corpus file, the analysis scripts, the power calculation
and the physics all regenerate from the repository, and every number in Chapters~5 and~6 is
produced by a script rather than transcribed.

\subsection{The framework}

The mean-field treatment of Chapter~3 derives its Landau coefficients from the
self-consistency condition rather than asserting them, recovers the exact critical exponents
numerically as a check on the solver, and locates the spinodal field $h_c(\beta J)$ at which
bistability collapses.

From that follows the bound: for any observable $X$ proposed as a field through $h = \alpha
X$, no media-driven transition is possible unless
\begin{equation}
\alpha \;\ge\; \frac{h_c(\beta J)}{\Delta X}.
\end{equation}
This inverts the usual procedure in the sociophysics literature, where the field is assigned
by hand and the population is then shown to tip. It requires no fit, applies to any candidate
observable, and can be evaluated before a corpus is collected.

\subsection{The measurement}

All 400 items were scanned and verified. The four categories separate at $\eta^2 = 0.1068$,
$p = 1.035\times10^{-9}$, against a design powered to detect $0.0268$. Against the labels the
items carried before scanning, the index achieves AUC $0.6274$, above the detectable floor of
$0.5916$ and in the predicted direction.

The separation is carried almost entirely by fear-activating content ($d = +0.939$). High
outrage, the category the proposal expected to separate most, does not separate at all
($d = +0.030$, $p = 0.832$).

Every category mean is negative. The deliberative-control network is predicted to respond
more strongly than the affective-salience network for all four categories; they differ in how
far below zero they sit, not in which network leads.

\section{What is claimed, and what is not}

\textbf{Claimed.} That the instrument works and is reproducible. That the index assigns
systematically different values to these four groups of documents, at an effect size the
design was powered to detect, with the power statement reported alongside. That the measured
spread $\Delta X = 0.1241$ places a lower bound on any coupling capable of driving a
transition, reaching $\alpha \ge 4.2938$ at $\beta J = 2$.

\textbf{Not claimed.} That the index detects manipulation. That the separation is due to
framing rather than to source dataset: category and provenance are confounded by construction
and this design cannot distinguish them. That the coupling has been measured: it is
unidentified, and no value of $\hat{\alpha}$ appears anywhere in this document. That anything
here measures a brain: every value is a prediction from a released encoder, the observable is
a cortical proxy, and no participant was scanned.

The gap between those two lists is the honest content of this work.

\section{Contribution}

Three things, in the order they are likely to be useful to someone else.

\textbf{The bound.} A necessary condition on the coupling between any content observable and
an opinion-dynamics field, requiring only the observable's measured spread and the model's own
spinodal. It is the part of this work that transfers beyond the instrument that produced it.

\textbf{The instrument, with its limitation documented.} A working, open pipeline from text to
predicted cortical response, together with the finding that the released checkpoint is
cortical-only, that the loader forces its average-subject configuration, and that its
agreement with real cortex is contested. A technical report describing an architecture is not
a warranty about released weights, and this project's experience is a concrete instance of
that.

\textbf{The measurement and its confound.} A separation reported with the design's
sensitivity attached and its central weakness stated before its interpretation, together with
the specific experiment that would resolve it.

\section{Future work}

\textbf{Remove the confound.} A corpus in which framing varies within a single source, or in
which sources are crossed with categories, would make the separation interpretable. It needs
no new instrumentation, only a rebuilt corpus, and it is the single change that most alters
what can be claimed.

\textbf{Finish the validation.} Whether the encoder tracks cortex at all is unresolved and
determines what the measurement means. The noise ceiling has been measured directly from
public data at $0.1517$, CI $[0.1484,\ 0.1551]$, and the evaluation code is written and
tested. What remains is running the checkpoint over a held-out stimulus and correlating, which
needs no large compute.

\textbf{A subcortical checkpoint.} The structures the original proposal named cannot be
measured with a cortical-only model. Training one means the full grid, the source fMRI
corpora, and compute outside the scope of an undergraduate project.

\section{Code, data and the working instrument}

The pipeline, the corpus builder, the analysis scripts and the figure scripts are public at
\url{https://github.com/brn-mwai/monarch}. Every number in this document is written by a
named script into a JSON or CSV artifact and read from there rather than transcribed, so any
figure or table can be regenerated from the repository alone.

The instrument itself runs at \url{https://monarch.brianmwai.com}. It carries the complete
400-item corpus with each item's text, its pre-scan labels and its measured values, and it
draws the predicted response on the cortical surface for the items whose per-vertex maps were
kept. It answers questions about the corpus from a generated fact sheet and declines anything
the data does not cover, including any request for a coupling value.

\section{Closing}

This work set out to make an invisible property of media measurable, and to say honestly what
the measurement supports. It delivers a working instrument, a bound that constrains any future
proposal of this kind, and a result that is real, modest, and confounded in a way this
document states plainly rather than leaves to be discovered.

The physics is the part that survives whatever becomes of the encoder. Whether or not the
released checkpoint predicts any brain, the requirement that $\alpha \ge h_c(\beta J)/\Delta
X$ holds for every observable anyone proposes as a field, and the burden it places on such
proposals does not depend on this instrument being the right one.

The longer aim behind the project, set out in \S\ref{sec:phase1}, is a tool that a reader
rather than a researcher could use to ask whether something they are about to watch is built
to manipulate them or to hold their attention. This document is the first phase of that and
makes no claim to be more. It leaves an instrument, a characterised measurement, an honest
account of where that measurement fails, and a bound that does not depend on any of the
above. What it does not leave is a tool anyone should use to judge a particular article or
programme, and the four reasons why are measured in these chapters rather than guessed at.
Stating the destination is not the same as arriving at it, and the difference between them is
most of what this thesis reports.
